\scanpage{103}

elle \deleted{vaut \unsure{à}} sur $1,U_1,U_2$ par

\[
1\longrightarrow2\deleted{\unsure{e}}\quad U_1\longrightarrow-b_1,\quad U_2\longmapsto-b_2
\]

\unsure{élémts|élémts.|étant} de $\Gamma\underline{\mathcal O}\subset\Gamma A$, il reste à trouver l'image de $U_1\otimes U_2$. C'est donc un \unsure{élémt|élément} de $A$ dont l'image dans $(A_1\otimes A_2)^{\unsure{\leftrightharpoons|\leftrightarrows|\natural}}$ est

\[
U_1\otimes U_2+\underbrace{\bar U_1\otimes\bar U_2}_{(-b_1\unsure{+|-}U_1)\otimes(-b_2-U_2)}
=\underbrace{2U_1U_2+b_1U_2+b_2U_1}_{\text{image de }U}+b_1b_2
\]

\deleted{donc} donc il y a lieu de prendre

\[
U_1*U_2=U+b_1b_2
\]

Mais il faut maintenant définir une application

\[
N_{A_1\otimes A_2/A}:A_1\otimes A_2\longrightarrow A
\]

application quadratique, dont la valeur en $z$ serait, essentiellement, $z\bar z$, mais ``remonté'' à $A$. L'application bil. associée serait

\[
(u,v)\longmapsto u\bar v+v\bar u=\mathrm{Tr}_N(u\bar v)\qquad\text{déjà connue}
\]

Donc il faut se donner les \deleted{\unsure{autre}} valeurs de $N$, et de $\varphi_N$, sur les élémts de la base $1,U_1,U_2,U_1\otimes U_2$ de $A_1\otimes A_2$. On \unsure{en|on} devrait avoir

\[
\begin{aligned}
N_A(x_1\otimes x_2)&=N_1(x_1)N(x_2)\in\Gamma\mathcal O\subset\Gamma A\\[-2pt]
\shortparallel &\\[-3pt]
(x_1\otimes x_2)(\bar x_1\otimes\bar x_2)&\\[-2pt]
\shortparallel &\\[-3pt]
(x_1\deleted{\unsure{\otimes}}\bar x_1)\otimes(x_2\bar x_{\unsure{1|2}})&
\end{aligned}
\]

En particulier

\[
\begin{aligned}
N(1)=1,\quad N(U_1)&=N_1(U_1)=c_1,\quad N(U_2)=N_2(U_2)=c_2\\
N(\deleted{\unsure{A}}U_1\otimes U_2)&=c_1c_2
\end{aligned}
\]

% archive sheet number: 102

